% Part: sets-functions-relations
% Chapter: sets
% Section: basics

\documentclass[../../../include/open-logic-section]{subfiles}

\begin{document}

\olfileid{sfr}{set}{bas}
\olsection{外延性}

一个\emph{集合}是对象的汇集，被视作单个对象。构成集合的对象称为该集合的\emph{元素}或\emph{成员}。若 $x$ 是集合~$A$ 的元素，则写作 $x \in A$；否则写作 $x \notin A$。没有元素的集合称为\emph{空集}，记作 “$\emptyset$”。

\begin{explain}
我们如何\emph{指定}集合，如何\emph{排列}其元素，乃至将其元素\emph{计数}多少次，都无关紧要。唯一重要的是它的元素是什么。我们将此概括为以下原理。
\end{explain}

\begin{defn}[外延性]
  若 $A$ 和 $B$ 是集合，则 $A = B$ 当且仅当 $A$ 的每个元素也都是 $B$ 的元素，反之亦然。
\end{defn}

外延性为某些记法提供了依据。一般来说，当我们有一些对象 $a_{1}$, \dots, $a_{n}$ 时，$\{a_{1}, \dots, a_{n}\}$ 就是\emph{那个}以 $a_1, \ldots, a_n$ 为元素的集合。我们强调“\emph{那个}”一词，因为外延性告诉我们这样的集合只能有\emph{一个}。事实上，外延性也为下式提供了依据：
  \[
    \{a, a, b\} = \{a, b\} = \{b,a\}.
  \]
这就说明了：当我们考虑集合时，我们并不关心其元素的顺序，也不关心它们被指定了多少次。

\begin{tagblock}{novice}
\begin{ex}
任何时候，只要有一堆对象，就可以把它们收集在一个集合里。例如，Richard（理查德）的兄弟姐妹构成的集合是一个只包含一个人的集合，我们可以把它写成 $S=\{\textrm{Ruth}\}$。小于 $4$ 的正整数集合是 $\{1, 2, 3\}$，但它也可以写成 $\{3, 2, 1\}$，甚至写成 $\{1, 2, 1, 2, 3\}$。根据外延性，这些都是同一个集合。因为 $\{1, 2, 3\}$ 的每个元素也都是 $\{3, 2, 1\}$（以及 $\{1,
2, 1, 2, 3\}$）的元素，反之亦然。
\end{ex}
\end{tagblock}

我们常常通过其元素共有的某个性质来指定一个集合。为此我们将使用以下简记：$\Setabs{x}{\phi(x)}$，其中 $\phi(x)$ 代表 $x$ 要成为该集合元素所必须具有的性质。

\begin{tagblock}{novice}
\begin{ex}
在我们的例子中，我们也可以将 $S$ 指定为
\[
S = \Setabs{x}{x \text{ 是 Richard 的兄弟姐妹}}.
\]
\end{ex}
\end{tagblock}

\begin{tagblock}{math}
\begin{ex}
一个数被称为\emph{完全数}，当且仅当它等于其真因子（即能整除它但不等于它自身的数）之和。例如，$6$ 是完全数，因为它的真因子是 $1$、$2$ 和~$3$，且 $6 = 1 + 2 + 3$。事实上，$6$ 是小于 $10$ 的正整数中唯一的完全数。因此，利用外延性，我们可以说：
\[
	\{6\} = \Setabs{x}{x\text{ 是完全数且 }0 \leq x \leq 10}
\]
我们将右边的记法读作“满足‘$x$ 是完全数且 $0 \leq x \leq 10$’的所有 $x$ 的集合”。此处的等式确认了，当我们考虑集合时，我们并不关心它们是如何被指定的。更一般地，外延性保证了使得 $\phi(x)$ 成立的 $x$ 的集合总是唯一的。因此，外延性为我们把 $\Setabs{x}{\phi(x)}$ 称为使得~$\phi(x)$ 成立的 $x$ 的\emph{那个}集合提供了依据。
\end{ex}
\end{tagblock}

外延性为我们提供了一种证明集合相等的方法：要证明 $A = B$，就要证明只要 $x \in A$ 则 $x \in B$，并且只要 $y \in B$ 则 $y \in A$。

\begin{prob}
证明最多只存在一个空集，即证明：如果 $A$ 和 $B$ 都是没有元素的集合，那么 $A = B$。
\end{prob}

\end{document}